
This study investigated whether P(True) calibration degrades on hard code problems for LLM code verifiers. The pre-registered hypothesis that DELTA\_ECE > 0 in at least 2/3 model families was refuted (1/3 pass). The actual finding is that DELTA\_ECE direction is architecture-determined: +0.298 for code-specialized (DeepSeek-Coder), -0.249 for code-adapted (CodeLlama), and approximately 0 for general-purpose (Llama3).

The methodology -- k=5 self-contained difficulty bootstrap, P(True) logprob extraction, M=15-bin ECE with bootstrap CI -- is validated and reusable. The infrastructure achieves perfect coverage across all model-benchmark combinations and produces non-degenerate confidence distributions for all tested models.

The practical implication is that P(True) calibration behavior is an architectural property, not a universal one. Code verification pipelines should characterize their model's difficulty-calibration profile before setting confidence thresholds.

Future directions include: (1) replication with k >= 20 for finer stratification and larger CodeLlama easy tier; (2) testing tier-specific temperature scaling (T\_hard not equal to T\_easy); (3) expanding to N >= 2 models per architecture category; (4) replication with instruction-tuned model variants to determine whether RLHF/SFT training alters the architecture-dependent DELTA\_ECE pattern; and (5) analysis of CodeLlama's training corpus overlap with MBPP easy-tier problems to test the overconfidence hypothesis.
